% =====================================================================
\section{Risk-Adaptive Routing: \mirage{} v5}\label{sec:mirage}
v1--v4 decide \emph{who} may call a tool and inspect \emph{what} crosses the
boundary, but every decision is binary (serve or refuse) and, at the content
layer, memoryless and threshold-keyed. Three gaps follow, and
\cref{sec:bastion} names all three: an adversary who stays just under the
behavioral thresholds is never flagged; a honeytoken yields attribution only
if the credential is \emph{reused} against a reporting service; and key/value
honeytokens change the value format, which can warn a sophisticated attacker.
\mirage{} (v5) replaces the binary post-authentication decision with a graded,
in-band one. It adds three coupled components and no external dependency
(standard library only).

\paragraph{A. Predictive session risk.}
For each authenticated \code{tools/call}, \mirage{} updates a per-source risk
state with a sequential change-point estimator over a feature stream
(inter-arrival time, argument entropy, tool diversity, schema-edge proximity).
A short warmup learns a per-feature mean and variance; the variance is then
\emph{frozen} as the in-control noise model, and the mean adapts only on
in-band requests ($|z|\le\code{Z\_SLACK}$), so an adversary cannot walk the
baseline along with a slow attack. A leaky one-sided CUSUM accrues only the
deviation \emph{beyond} the per-feature noise band, so a run of
individually-permissible requests that collectively drift crosses the boundary
(the low-and-slow case) while isolated spikes decay. No fixed per-request count
is used.

\paragraph{B. Graded execution router.}
The risk state selects a lane per request: \code{PRIMARY} (production
resources), \code{THROTTLED}, \code{SHADOW}, or \code{DECEPTION}. The shadow
lane executes a \emph{functional replica of the same tool} against a
production-isolated synthetic data environment, returning a schema-conforming
result derived from synthetic data. Because the decoy is the tool's own code
over synthetic data, its output is format- and distribution-matched (closing
the tip-off gap) while touching no real resource. Routing is transparent and
sticky: once a session crosses into the shadow lane it stays there, so an
adversary cannot probe back to the real lane by briefly behaving. The deception
lane additionally seeds the synthetic result with a canary whose identifier is
bound into the provenance record, so a silent collector that never reuses the
credential still interacts with tracked synthetic data.

\paragraph{C. Verifiable result provenance.}
Each returned result is bound to a non-forgeable identifier of the lane that
produced it and to a running session digest chaining it to all prior results
($\text{chain}_i = \mathrm{HMAC}(k,\ \text{chain}_{i-1}\,\|\,i\,\|\,
\text{lane}_i\,\|\,\mathrm{H}(\text{result}_i))$). A consuming agent verifies
the chain \emph{before} admitting a result into the model's context, detecting
tampering, lane forgery, reordering, or out-of-session injection.

\finding{Deployment and evaluation.}
\mirage{} runs as a loopback-bound service on the study host (it acts only on
\emph{authenticated} calls, so no public exposure is required or meaningful).
Over a clean window the service logged \num{\MrgTotal} authenticated requests
across benign and adversarial sessions (\cref{tab:mirage}; aggregates from the
live \code{evidence\_report.json}). On \num{\MrgBenignN} benign requests with
realistic per-request jitter, \mirage{} shadowed \MrgBenignShadow{}
(\MrgBenignRate\%), preserving the v1--v4 property of no false action on benign
sessions. A modeled low-and-slow adversary (a steady sub-threshold drift) was
contained to the shadow lane after \MrgSlowFirst{} calls and never again served
production data; an overtly anomalous session was isolated on its first
post-warmup call with \MrgOvertLeak{} real-data responses after containment.
Over the whole window the lane mix was \num{\MrgLanePrimary} primary,
\MrgLaneThrottled{} throttled, \MrgLaneShadow{} shadow and \MrgLaneDeception{}
deception.

\finding{Cost.}
The added hot path (risk update, routing, attestation) is
\SI{\MrgOverheadUs}{\micro\second} per request in pure-Python standard library,
the same order as the \sys{} gate's \SI{9}{\micro\second}. Verifying a full
\num{\MrgVerifyChain}-call session chain costs \SI{\MrgVerifyMs}{\milli\second}.
Source tracking reuses the \sys{} LRU bound (\num{20000} sources).

\begin{table}[t]
\centering\small
\caption{\mirage{} v5 on the deployed loopback service
(\code{evidence\_report.json}; \num{\MrgTotal} authenticated requests).}
\label{tab:mirage}
\begin{tabularx}{\columnwidth}{@{}lX@{}}
\toprule
\# & Result \\
\midrule
M1 & Benign false-shadow: \textbf{\MrgBenignShadow/\num{\MrgBenignN}}
     (\MrgBenignRate\%). \\
M2 & Low-and-slow drift (\num{\MrgSlowN} calls): contained to shadow after
     \MrgSlowFirst{} calls; \MrgOvertLeak{} production responses thereafter. \\
M3 & Overt anomaly (\num{\MrgOvertN} calls): shadowed on the first post-warmup
     call; \textbf{\MrgOvertLeak} real-data leaks after containment. \\
M4 & Provenance: a session chain verifies clean; a single swapped result is
     caught at its exact sequence index. \\
\midrule
E1 & Added overhead \SI{\MrgOverheadUs}{\micro\second}/request (stdlib). \\
E2 & Verify \num{\MrgVerifyChain}-call chain: \SI{\MrgVerifyMs}{\milli\second}. \\
\bottomrule
\end{tabularx}
\end{table}

\finding{Limitations.}
The estimator needs a short warmup before the CUSUM is trusted, so the first
few calls of a session are served from the primary lane; a slow adversary
receives real data until the drift is detected (\MrgSlowFirst{} calls in M2).
The \MrgBenignRate\% false-shadow figure is measured against a specific benign
jitter model; like the v2--v4 thresholds, \mirage{}'s \code{Z\_SLACK},
\code{CUSUM\_LEAK} and alarm boundary are reasonable defaults, not tuned
optima, and trade detection latency against false-shadow rate. Provenance binds
\emph{which lane} produced a result and detects tampering, but does not by
itself establish that the primary lane's data is correct.
